\section{Implementation}

The code is a Python package (\texttt{cpt\_budget}) whose modules mirror the information boundary:
\begin{itemize}
\item data splits;
\item models with a common \texttt{condition(revealed)}/\texttt{predict(query)} interface, which strips
  queries of measured values;
\item acquisition and stop rule;
\item replay loop;
\item a separate audit evaluator;
\item economics.
\end{itemize}
The frozen configuration stores SHA-256 fingerprints of data, weights and settings. The core of batch
selection is:
\begin{lstlisting}
while remaining and chosen_count < batch_size:
    dist  = [min_distance(g, observed_xy) for g in remaining]
    score = 0.5*group.rank_U + 0.3*group.rank_D + 0.2*rank01(dist)
    gid   = argmax(score)                         # ties: group id, then location id
    chosen += members(gid)                        # whole near-duplicate group
    observed_xy = vstack([observed_xy, xy(gid)])  # coordinates only, never curves
\end{lstlisting}
Each round is logged; for example, the adaptive run at Wellington Airport (seed 17) stops in round 3:
\begin{lstlisting}[style=output]
round n_obs dQ90  dH90 support_logMAE gp_std_P90 streak stop
  3    11   0.067 0.00     0.132        0.169      2    online_stability
\end{lstlisting}
\textbf{Key decisions.} TabPFN weights stay frozen, so each round only updates the context and takes
seconds. Grouping is enforced throughout. All 14 CNP and 129 TabPFN weight tensors differ from
initialisation, so weak CNP results are not an untrained-model artefact.

\textbf{Checks.} Automated checks show that audit labels cannot affect predictions (24 paired contexts,
difference 0). They also confirm that group counts are exact, inputs are unchanged and economic sums
recompute. All 95 unit tests pass.
